\documentclass{beamer}
\usepackage[utf8]{inputenc}
\usepackage{amsmath}

\usetheme{AnnArbor}

\title{Curvatura y Torsión, Vector Tangente Unitario y Recta Tangente}
\author{David E. Duarte G.}
\date{}

\begin{document}
	
	\frame{\titlepage}
	
	\begin{frame}{Introducción}
		En esta presentación estudiaremos conceptos fundamentales de la geometría diferencial, como la curvatura y torsión de una curva en el espacio tridimensional, así como el vector tangente unitario y la recta tangente en un punto dado.
	\end{frame}
	
	\begin{frame}{Curvatura y Torsión}
		\begin{itemize}
			\item La curvatura de una curva en el espacio se define como el cambio en la dirección del vector tangente unitario en función de la longitud de arco.
			\item La torsión mide la variación en la orientación del plano osculador de la curva.
		\end{itemize}
	\end{frame}
	
	\begin{frame}{Fórmulas}
		\begin{align*}
			\kappa &= \frac{\lvert \textbf{r}'(s) \times \textbf{r}''(s) \rvert}{\lvert \textbf{r}'(s) \rvert^3} \\
			\tau &= \frac{\textbf{r}'(s) \cdot (\textbf{r}''(s) \times \textbf{r}'''(s))}{\lvert \textbf{r}'(s) \times \textbf{r}''(s) \rvert^2}
		\end{align*}
	\end{frame}
	
	\begin{frame}{Vector Tangente Unitario}
		\begin{itemize}
			\item El vector tangente unitario a una curva en un punto dado se define como la derivada del vector posición con respecto a la longitud de arco.
		\end{itemize}
	\end{frame}
	
	\begin{frame}{Fórmula}
		\begin{equation*}
			\textbf{T}(s) = \frac{\textbf{r}'(s)}{\lvert \textbf{r}'(s) \rvert}
		\end{equation*}
	\end{frame}
	
	\begin{frame}{Recta Tangente}
		\begin{itemize}
			\item La recta tangente a una curva en un punto dado es la recta que pasa por dicho punto y tiene la misma dirección que el vector tangente unitario en ese punto.
		\end{itemize}
	\end{frame}
	
	\begin{frame}{Ecuación}
		La ecuación de la recta tangente en el punto $P$ de la curva $\textbf{r}(s)$ es:
		\begin{equation*}
			\textbf{r}(s) = \textbf{r}(s_0) + (\textbf{r}'(s_0))(s - s_0)
		\end{equation*}
	\end{frame}
	
	\begin{frame}{Conclusiones}
		Hemos revisado conceptos fundamentales de la geometría diferencial, como la curvatura y torsión de una curva en el espacio tridimensional, así como el vector tangente unitario y la recta tangente en un punto dado.
	\end{frame}
	
\end{document}
